% =====================================================================
%  Secao 6: O Processo de Retrocesso
% =====================================================================
\section{O Processo de Retrocesso}

O nome Backtracking vem justamente da capacidade de retornar para decisões anteriores.

Suponha que um algoritmo esteja tentando resolver um problema. Após realizar várias escolhas consecutivas, ele descobre que a última decisão torna impossível alcançar uma solução.

Nesse momento, em vez de continuar desperdiçando recursos, o algoritmo volta para o último ponto de decisão e tenta outra alternativa. Esse processo pode ocorrer diversas vezes até que uma solução válida seja encontrada.

\begin{passos}[Em outras palavras]
Escolhe. Testa. Avança. Erra. Volta. Tenta novamente.
\end{passos}

\figura{fig:retrocesso}%
  {Ao encontrar um caminho inviável, o algoritmo retorna ao último ponto de decisão e tenta outra alternativa.}%
  {figuras/tikz/retrocesso}

\begin{ideiachave}
Esse comportamento é o coração do Backtracking.
\end{ideiachave}
